\chapter{Endoscopic Sinus Surgery}

\section*{Definition and Overview}
Endoscopic sinus surgery is a procedure used to treat chronic sinusitis that has not improved with medical treatment. A thin, rigid telescope called an endoscope is passed through the nostrils, allowing the surgeon to see inside the sinuses without making any cuts on the face. Using fine instruments, the surgeon opens the blocked sinus openings, removes diseased tissue and polyps, and restores drainage. The surgery is performed under general anaesthetic and is usually done as a day procedure, with most people going home the same day.

\section*{Epidemiology}
Chronic sinusitis is one of the most common chronic conditions in Australia, affecting a significant proportion of the population. Most people are managed successfully with medicines, but a minority have persistent disease despite treatment and benefit from surgery. Endoscopic sinus surgery is among the most common ear, nose, and throat operations, and its use has increased as the technique has become refined and less invasive.

\section*{Aetiology and Pathophysiology}
Sinusitis is inflammation of the lining of the sinuses, the air-filled spaces around the nose.
\begin{itemize}
    \item \textbf{Blockage:} the narrow openings that drain the sinuses become blocked, preventing drainage and ventilation
    \item \textbf{Inflammation:} infection, allergy, and structural factors cause persistent inflammation of the lining
    \item \textbf{Polyps:} nasal polyps can block the sinus openings and are a common reason for surgery
    \item \textbf{Chronic infection:} bacteria and fungi can persist in sinuses that do not drain
    \item \textbf{The operation:} enlarging the sinus openings, removing polyps and diseased tissue, and restoring normal drainage
    \item \textbf{Why surgery:} it removes the blockage and allows medical treatment to work more effectively
\end{itemize}

\section*{Clinical Features}
Surgery is considered for chronic sinusitis with particular features.
\begin{itemize}
    \item Persistent nasal blockage, facial pain, and pressure
    \item Thick nasal discharge and reduced sense of smell
    \item Headaches and tiredness
    \item Recurrent acute sinus infections
    \item Nasal polyps
    \item Symptoms that have not responded to medical treatment
    \item Sinus disease affecting the eyes or causing complications
\end{itemize}

\section*{Differential Diagnosis}
Before surgery, the diagnosis and the reasons for failure of medical treatment are reviewed.
\begin{itemize}
    \item Allergic rhinitis and other causes of nasal blockage
    \item Migraine and other causes of facial pain
    \item Dental causes of facial pain
    \item Structural problems, including a deviated septum
    \item Fungal sinusitis, which may need specific management
    \item Immune deficiency or cystic fibrosis in recurrent disease
\end{itemize}

\section*{When to Seek Medical Care}
People with chronic sinusitis should be assessed by a GP and referred to an ear, nose, and throat specialist when medical treatment fails. Urgent care is needed for complications such as eye swelling or severe headache.

\textbf{Call 000 (or local emergency services) for:}
\begin{itemize}
    \item Swelling, redness, or pain around the eye, which may indicate spread of infection
    \item Severe headache with fever and confusion
    \item Sudden vision changes
    \item Stiff neck or signs of meningitis
\end{itemize}

\section*{Diagnosis and Investigations}
Before surgery, the sinuses are fully assessed.
\begin{itemize}
    \item \textbf{Nasal endoscopy:} examination of the nose and sinus openings with a small camera
    \item \textbf{CT scan:} the key imaging test, showing the extent of disease and the anatomy
    \item \textbf{Allergy testing:} for allergic causes
    \item \textbf{Swabs:} for infection in selected cases
    \item \textbf{Pre-operative assessment:} general health and anaesthetic review
\end{itemize}

\section*{Management}
The operation and recovery are managed in a structured way.
\begin{itemize}
    \item \textbf{The procedure:} under general anaesthetic, the surgeon uses an endoscope to open the sinuses and remove disease
    \item \textbf{Minimally invasive:} no external cuts, and in some cases balloon sinuplasty is used
    \item \textbf{Day surgery:} most people go home the same day
    \item \textbf{Aftercare:} nasal saline rinses, and sometimes packing or dressings
    \item \textbf{Medical treatment:} nasal sprays, and antibiotics or steroids for a period after surgery
    \item \textbf{Follow-up:} endoscopic review to check healing
    \item \textbf{Continued care:} ongoing medical treatment is often needed, as surgery controls but does not always cure sinus disease
\end{itemize}

\section*{Complications}
\begin{itemize}
    \item Bleeding, which is usually minor but can rarely be significant
    \item Infection
    \item Scarring that narrows the openings again
    \item Damage to nearby structures, including the eye and the base of the skull, which is rare
    \item Reduced sense of smell, which may persist
    \item Recurrence of polyps or sinusitis
    \item The effects of general anaesthesia
\end{itemize}

\section*{Prognosis and Outcomes}
Endoscopic sinus surgery is effective for most people with chronic sinusitis, improving nasal breathing, reducing facial pain, and cutting the frequency of infections. People with polyps usually improve markedly, though long-term medical treatment is often needed to prevent recurrence. Overall satisfaction is high, and serious complications are rare.

\section*{Prevention and Self-Care}
\begin{itemize}
    \item Use nasal saline rinses after surgery as advised
    \item Continue prescribed nasal sprays even when feeling well
    \item Avoid irritants and allergens that trigger sinusitis
    \item Attend follow-up appointments for endoscopic review
    \item Seek prompt treatment of infections
    \item Do not smoke, as it worsens sinus disease
    \item Manage underlying conditions such as allergy and reflux
\end{itemize}

\section*{Sources and Further Reading}
The following reputable sources provide further information for healthcare students and patients seeking reliable, current advice about endoscopic sinus surgery.
\begin{itemize}
    \item Australian Society of Otolaryngology Head and Neck Surgery. \textit{Sinus surgery information.} https://www.asohns.org.au/
    \item Healthdirect Australia. \textit{Sinusitis and sinus surgery.} https://www.healthdirect.gov.au/sinusitis
    \item NHS. \textit{Sinus surgery (functional endoscopic sinus surgery).} https://www.nhs.uk/conditions/sinusitis/
    \item Mayo Clinic. \textit{Nasal and sinus surgery.} https://www.mayoclinic.org/diseases-conditions/chronic-sinusitis/
    \item MedlinePlus. \textit{Sinus surgery.} https://medlineplus.gov/ency/article/007621.htm
    \item Better Health Channel. \textit{Sinusitis.} https://www.betterhealth.vic.gov.au/
\end{itemize}
